% !PRINTED 615
\weberpage{615}
\gesperrt{Bedeutet $E$ die Gruppe der Einheiten oder auch eine darin enthaltene Gruppe, so bekommen wir denselben Wert der Klassenzahl und denselben Klassenkörper, wenn wir $EA$ an Stelle von $A$ treten lassen.}

Dies gilt nicht nur für diese besondere Gruppe, sondern allgemein. Wir hätten also z.\,B. die Zahlen $\alpha$ der Gruppe $A$ auch durch
\[
\alpha \equiv \pm 1 \quad (\mathrm{mod}\ \mathfrak{m})
\]
definieren können, denn $E = +1, -1$ ist auch in den Ausnahmefällen $\Delta = -4$, $\Delta = -3$ eine Gruppe.

Nun läßt sich nachweisen, daß der Teilungskörper $\mathfrak{T}_{\mathfrak{m}}$ der Klassenkörper zu $A$ ist. Schließen wir zunächst noch alle in der Diskriminante des Teilungskörpers aufgehenden Primfaktoren aus, und bezeichnen mit $\tau$ irgend eine Zahl des Teilungskörpers, so ist nach \S\,160
\weq{5}{\tau^p \equiv \tau \quad (\mathrm{mod}\ \mathfrak{p})}
nur dann erfüllt, wenn $\mathfrak{p}$ ein \gesperrt{Primideal ersten Grades} in $\Omega$ ist, das der Hauptklasse angehört. Das aber ist das Kennzeichen des Klassenkörpers. Der Grad des Teilungskörpers ist aber nach \S\,154 höchstens gleich dem in (4) gegebenen Ausdruck $h$, folglich ist er genau gleich $h$ und die Teilungsgleichung irreduzibel. Außerdem ist damit bewiesen, daß es in jeder der Idealklassen $\overline{A}_\nu$ nach $A$ unendlich viele Primideale ersten Grades gibt.

Man kann diesem Satz folgenden Ausdruck geben:

Ist $\beta$ eine beliebige Zahl in $O$, $\mathfrak{a}$ ein beliebiges Ideal in $\overline{O}$, so ist durch
\weq{6}{\overline{A}_\nu = \mathfrak{a}\beta A}
eine Klasse nach $A$ bestimmt (\S\,161), und in dieser Klasse sind unendlich viele Primideale ersten Grades enthalten. Ein solches Primideal ist nach (6) in der Form darstellbar:
\weq{7}{\mathfrak{p} = \mathfrak{a}\beta\alpha,}
worin $\alpha \equiv 1$ (mod $\mathfrak{m}$) ist. Setzen wir also
\[
\alpha\beta = \pi,
\]
so ist
\[
\pi \equiv \beta \quad (\mathrm{mod}\ \mathfrak{m}).
\]
Daraus folgt:

2. \gesperrt{Ist $\beta$ eine Zahl in $O$, $\mathfrak{a}$ ein Ideal in $\overline{O}$, so gibt es unendlich viele im allgemeinen gebrochene Zahlen $\pi$ in $O$, die der Bedingung}
%CONT>>
